\section{Сборка и проверка работы}
\subsection{Среда сборки}
Проверка выполнена на одном компьютере в двух средах.

\begin{center}
\begin{tabular}{p{0.2\textwidth}p{0.7\textwidth}}
Среда & Состав \\
\hline
Windows & Microsoft Windows NT 10.0.26200.0, Visual Studio 18 2026, компилятор MSVC 19.51.36260.0, конфигурация Release \\
Linux & WSL2, Ubuntu 24.04.3 LTS, ядро 6.6.87.2-microsoft-standard-WSL2, g++ 13.3.0, CMake 3.28.3, Python 3.12.3, конфигурация Release \\
\end{tabular}
\end{center}

Из корня проекта \texttt{Operation-Systems} для Windows выполнены команды:
\begin{verbatim}
cmake -S . -B build-lab4
cmake --build build-lab4 --config Release --target lab4_parent lab4_child
\end{verbatim}

Для Linux:
\begin{verbatim}
cmake -S . -B /tmp/lab4-wsl -DCMAKE_BUILD_TYPE=Release
cmake --build /tmp/lab4-wsl --target lab4_parent lab4_child
\end{verbatim}

Сборка в обеих средах завершилась без ошибок компиляции и компоновки. В Windows созданы файлы \texttt{lab4\_parent.exe} и \texttt{lab4\_child.exe}, в Linux --- \texttt{lab4\_parent} и \texttt{lab4\_child}. Обе программы должны находиться в одной папке: родитель ищет ребёнка рядом с собой.

\subsection{Контрольный пример}
Последовательность ввода:
\begin{verbatim}
result.txt
4
6
9
17
\end{verbatim}

Числа 4, 6 и 9 составные и записываются в файл, число 17 простое и завершает работу обоих процессов. Ниже приведён фактический вывод программы в Linux и содержимое выходного файла. Приглашения идут подряд, так как ввод подавался из файла.

\lstinputlisting[language={},basicstyle=\footnotesize\ttfamily,numbers=none,breaklines=true,caption={Запуск в Linux (WSL2)},captionpos=b]{\labroot/inc/run_linux.txt}

Тот же сценарий в Windows дал тот же вывод и то же содержимое файла; после завершения программы во временной папке не осталось файлов \texttt{lm4*.tmp}.

\lstinputlisting[language={},basicstyle=\footnotesize\ttfamily,numbers=none,breaklines=true,caption={Запуск в Windows},captionpos=b]{\labroot/inc/run_windows.txt}

\subsection{Трассировка системных вызовов}
Для подтверждения того, что обе программы используют один отображаемый файл, запуск в Linux был проверен командой \texttt{strace -f -e trace=execve,openat,ftruncate,mmap,unlink,clone}. Из журнала отобраны строки, относящиеся к отображаемому файлу.

\lstinputlisting[language={},basicstyle=\scriptsize\ttfamily,numbers=none,breaklines=true,caption={Фрагмент трассировки (процесс 984 --- родитель, 985 --- ребёнок)},captionpos=b]{\labroot/inc/strace_lab4.txt}

В журнале видно следующее. Родитель создаёт файл \texttt{/tmp/lab4\_map\_y4gIiw} и отображает 600 байт с флагом \texttt{MAP\_SHARED}: это общая структура и два семафора. Затем создаётся процесс 985, который запускает \texttt{lab4\_child} с путём к файлу первым аргументом. Ребёнок открывает тот же файл и отображает тот же размер. В конце родитель удаляет файл вызовом \texttt{unlink}~\cite{linuxmmap}.

\subsection{Автоматические тесты}
Из корня проекта запускается скрипт, который создаёт настоящие процессы во временных папках:
\begin{verbatim}
python lab4/tests/test_lab4.py build-lab4/bin/lab4/lab4_parent.exe
python3 lab4/tests/test_lab4.py /tmp/lab4-wsl/bin/lab4/lab4_parent
\end{verbatim}

В Windows пройдено 25 тестов, в Linux --- 26: дополнительный тест в Linux проверяет, что при уничтожении ребёнка сигналом \texttt{SIGKILL} родитель сообщает об ошибке и не зависает. В Windows этот тест пропускается, так как для поиска процесса используется утилита \texttt{pgrep}.

\begin{center}
\begin{tabular}{p{0.34\textwidth}p{0.56\textwidth}}
Входные данные & Ожидаемый результат \\
\hline
\texttt{4, 6, 9, 17, 25} & Запись 4, 6, 9, завершение на 17 \\
\texttt{8, -4, 10} & Запись 8, завершение на отрицательном значении \\
\texttt{2} & Завершение без записи \\
\texttt{0, 1, 4, -1} & Пропуск 0 и 1, запись 4 \\
Конец ввода & Команда \texttt{Stop}, нормальное завершение \\
Текст, два числа, пустая строка & Сообщение об ошибке ввода, обмен продолжается \\
Значения вне \texttt{int} & Отклонение команды \\
2147483647 и 2147483646 & Простое число завершает работу, чётное записывается \\
Строка 200000 символов & Отклонение: строка не помещается в буфер \\
Строка 251 и 256 байт & Первая принимается, вторая отклоняется \\
Имя файла с пробелами и русскими буквами & Файл создан и заполнен \\
Недоступный выходной каталог & Ошибка и код завершения 1 \\
Отсутствует \texttt{lab4\_child} & Ошибка запуска, ненулевой код \\
2000 команд подряд & Все составные числа записаны в том же порядке \\
Уничтоженный ребёнок (Linux) & Сообщение об ошибке, родитель не зависает \\
\end{tabular}
\end{center}

После всех тестов во временной папке не остаётся файлов отображения. Измерения производительности не проводились: они не входят в задание. Проверка Linux-реализации выполнена в WSL2, а не на отдельном компьютере с Linux.
